%=============================================================================
% ENGINEERING DESIGN: TERRESTRIAL psi-PNT DEMONSTRATION NETWORK
%
% Patent #9: psi-Field Navigation, Positioning, and Timing Systems
% Author: Gary Thomas Alcock
% Density Field Dynamics Research Programme
%=============================================================================
\documentclass[11pt,a4paper]{article}

\usepackage{amsmath,amssymb}
\usepackage{siunitx}
\usepackage{graphicx}
\usepackage{booktabs}
\usepackage{hyperref}
\usepackage{geometry}
\usepackage{longtable}
\usepackage{enumitem}
\usepackage{xcolor}
\usepackage{fancyhdr}

\geometry{margin=2.5cm}
\pagestyle{fancy}
\fancyhf{}
\fancyhead[L]{\small $\psi$-PNT Demonstration Network -- Engineering Design}
\fancyhead[R]{\small DFD Patent \#9}
\fancyfoot[C]{\thepage}

\title{Engineering Design Document:\\
Terrestrial $\psi$-PNT Demonstration Network\\[0.5em]
\large A GPS-Free, Spoof-Immune Positioning System\\
Based on Scalar Refractive Field Geometry}

\author{Gary Thomas Alcock\\
\small Density Field Dynamics Research Programme}

\date{Document Reference: DFD-P09-ENG-001 \quad Rev.~A\\[0.3em] \today}

\begin{document}
\maketitle
\tableofcontents
\newpage

%=============================================================================
\section{Executive Summary}
%=============================================================================

This document presents the complete engineering design for a terrestrial
demonstration network of the $\psi$-field Positioning, Navigation, and
Timing ($\psi$-PNT) system described in DFD Patent \#9. The network
comprises six emitter stations, a central timing hub, mobile receiver
units, and a reference GNSS system for performance validation. The design
covers site selection, emitter engineering, receiver architecture, timing
distribution, communications, power systems, and a phased test program.

\paragraph{Design objectives.}
\begin{enumerate}[label=(\alph*)]
\item Demonstrate detection of modulated $\psi$-field signals from
  engineered emitter assemblies at ranges up to 5~km.
\item Demonstrate position determination using $\psi$-field
  measurements from $\geq 3$ emitters.
\item Demonstrate clock-free timing extraction from $\psi$-oscillation
  phase.
\item Demonstrate anti-spoofing via curvature tensor verification.
\item Demonstrate hybrid GNSS+$\psi$ fusion with spoofing detection.
\item Characterize system performance against Cram\'er-Rao bounds.
\end{enumerate}

%=============================================================================
\section{System Overview}
%=============================================================================

\subsection{Network Topology}

The demonstration network deploys six $\psi$-emitter stations in a
regular hexagonal arrangement with 5~km inter-station spacing,
providing coverage over a $\sim$100~km$^2$ test area. A central timing
hub at the geometric center distributes phase-coherent timing to all
stations.

\begin{center}
\begin{tabular}{cccc}
\toprule
\textbf{Station} & \textbf{Designation} & \textbf{Position (km)} & \textbf{Mod.\ freq.\ (Hz)} \\
\midrule
E1 & North & $(0.00, +5.00)$ & 10.00 \\
E2 & Northeast & $(+4.33, +2.50)$ & 12.50 \\
E3 & Southeast & $(+4.33, -2.50)$ & 15.00 \\
E4 & South & $(0.00, -5.00)$ & 17.50 \\
E5 & Southwest & $(-4.33, -2.50)$ & 20.00 \\
E6 & Northwest & $(-4.33, +2.50)$ & 22.50 \\
Hub & Central & $(0.00, 0.00)$ & --- \\
\bottomrule
\end{tabular}
\end{center}

Modulation frequencies are chosen with 2.5~Hz spacing to ensure
spectral separation at the receiver. The frequency range 10--22.5~Hz
lies within the optimal band for mechanical rotating mass assemblies
and avoids seismic noise peaks (microseismic band 0.1--0.3~Hz,
cultural noise $>30$~Hz).

\subsection{System Block Diagram}

The system comprises four principal subsystems:
\begin{enumerate}
\item \textbf{Emitter Subsystem:} Six identical emitter stations, each
  generating a modulated $\psi$-field signal.
\item \textbf{Timing Distribution Subsystem:} Central hub with master
  clock and phase-stabilized fiber links to each station.
\item \textbf{Receiver Subsystem:} Mobile and fixed receiver units
  measuring $\psi$, $\nabla\psi$, and extracting PVT solutions.
\item \textbf{Reference Subsystem:} GNSS/RTK network providing
  ground-truth positioning for performance validation.
\end{enumerate}

%=============================================================================
\section{Emitter Station Design}
%=============================================================================

\subsection{Modulated Mass Assembly}

Each emitter generates a time-varying $\psi$-field using a rapidly
rotating high-density mass. The design parameters are:

\subsubsection{Rotor Design}

\begin{longtable}{p{5cm}p{8cm}}
\toprule
\textbf{Parameter} & \textbf{Specification} \\
\midrule
\endfirsthead
\midrule
\textbf{Parameter} & \textbf{Specification} \\
\midrule
\endhead
Material & Tungsten alloy (W-Ni-Fe, $\rho = 18{,}500$~kg/m$^3$) \\
Rotor geometry & Dumbbell (two lobes on shaft) \\
Total rotor mass & 500~kg per station \\
Lobe mass & 200~kg each (400~kg total in lobes) \\
Shaft/hub mass & 100~kg \\
Lobe dimensions & Cylindrical, $r = 12$~cm, $L = 15$~cm \\
Lobe center-to-center & 60~cm \\
Rotation frequency & Station-specific: 5.00--11.25~Hz \\
$\psi$-modulation frequency & $2\times$ rotation (quadrupole): 10.00--22.50~Hz \\
Tip speed & $\leq 42$~m/s (well below burst limit) \\
Balance specification & $< 10$~g$\cdot$mm residual imbalance \\
Bearing type & Active magnetic bearings (AMB) \\
Bearing capacity & 10~kN radial, 5~kN axial \\
Vacuum enclosure & $< 10^{-3}$~Pa (minimize aerodynamic drag) \\
Drive motor & Permanent magnet synchronous, 5~kW \\
Speed control & Phase-locked to timing reference, $\delta f/f < 10^{-10}$ \\
\bottomrule
\end{longtable}

\subsubsection{$\psi$-Field Signal Calculation}

A dumbbell rotor with two lobes of mass $m_L$ at distance $d/2$ from
the rotation axis generates a quadrupole $\psi$-field. At distance $R$
from the rotor center ($R \gg d$), the time-varying component is:
\begin{equation}
\delta\psi(R, \theta, t) = \frac{Gm_L d^2}{4c^2 R^3}
\sin^2\theta\,\cos(2\omega_{\rm rot} t + 2\phi_0),
\end{equation}
where $\theta$ is the polar angle from the rotation axis and $\phi_0$
is the initial phase. At the equatorial plane ($\theta = \pi/2$) and
$R = 5$~km:
\begin{align}
\delta\psi(5\;\text{km}) &= \frac{6.67\times10^{-11}\times 200 \times 0.36}
{4 \times (3\times10^8)^2 \times (5\times10^3)^3} \nonumber\\
&= \frac{4.80\times10^{-9}}{4.5\times10^{28}} \nonumber\\
&\approx 1.1\times10^{-37}.
\end{align}

This signal is extremely small with passive gravitational coupling alone.
However, in the DFD framework, the $\psi$-emitter couples to the scalar
field through the engineered refractive-index modulation, not merely through
Newtonian gravitational coupling. The DFD coupling enhances the effective
signal by a factor related to the $\psi$-field amplification mechanisms
discussed in the companion paper.

\subsubsection{DFD-Enhanced Emission}

In DFD, the scalar field $\psi$ couples to matter density through the
action
\begin{equation}
S_{\rm coupling} = -\frac{c^2}{2}\int\psi(\rho - \bar{\rho})\,dt\,d^3x.
\end{equation}
An oscillating mass density $\rho(\mathbf{x},t)$ drives $\psi$-oscillations
through the field equation. The key enhancement over passive gravitational
coupling arises from resonant amplification when the driving frequency
matches a natural mode of the local $\psi$-field configuration.

For the demonstration network, the emitter design targets an effective
modulated signal at the receiver of:
\begin{equation}
\psi_{\rm eff}(R) = \frac{G_{\rm eff} M_E}{c^2 R},
\end{equation}
where $G_{\rm eff}$ incorporates the DFD coupling enhancement. The
required $G_{\rm eff}/G$ ratio for a detectable signal
($\psi_{\rm eff} \sim 10^{-20}$ at $R = 1$~km) is:
\begin{equation}
\frac{G_{\rm eff}}{G} \sim \frac{\psi_{\rm eff}\,c^2 R}{G M_E}
= \frac{10^{-20}\times 9\times10^{16}\times 10^3}{6.67\times10^{-11}\times 500}
\sim 2.7\times10^{7}.
\end{equation}

The demonstration network is designed to characterize the actual
DFD coupling enhancement and establish whether the theoretical
amplification mechanisms produce detectable signals.

\subsection{Station Infrastructure}

\subsubsection{Building and Foundation}

\begin{longtable}{p{5cm}p{8cm}}
\toprule
\textbf{Component} & \textbf{Specification} \\
\midrule
\endfirsthead
Building type & Reinforced concrete shelter, $5\times5\times3$~m \\
Foundation & Isolated mass block, 20~tonnes, vibration-decoupled \\
Floor isolation & Air spring mounts, $f_0 < 0.5$~Hz \\
Temperature control & HVAC, $\pm 0.5$~K stability \\
Humidity control & $< 50\%$ RH, dehumidifier \\
EMI shielding & Faraday cage (copper mesh), $> 60$~dB at 1~GHz \\
Security & Fenced compound, access control, CCTV \\
\bottomrule
\end{longtable}

\subsubsection{Power System}

\begin{longtable}{p{5cm}p{8cm}}
\toprule
\textbf{Component} & \textbf{Specification} \\
\midrule
\endfirsthead
Primary power & Grid connection, 3-phase 400~V, 50~kVA \\
UPS & Online double-conversion, 20~kVA, 30~min \\
Backup generator & Diesel, 30~kVA, auto-start \\
Rotor drive & 5~kW continuous + 2~kW magnetic bearings \\
Environmental systems & 5~kW HVAC + 2~kW monitoring \\
Timing/comms & 1~kW \\
Total station load & $\sim 15$~kW continuous \\
\bottomrule
\end{longtable}

\subsubsection{Environmental Monitoring}

Each station is equipped with:
\begin{itemize}
\item Broadband seismometer (STS-2 class, 0.01--50~Hz, 170~dB dynamic range)
\item Tiltmeter (resolution $< 1$~nrad)
\item Temperature sensors (Pt100, $\pm 0.01$~K accuracy) at rotor, floor, walls
\item Barometric pressure sensor ($\pm 0.1$~hPa)
\item Magnetic field sensor (fluxgate, 3-axis, $\pm 0.1$~nT)
\item Groundwater level sensor (if applicable)
\item GNSS antenna for station position monitoring
\end{itemize}

%=============================================================================
\section{Timing Distribution Subsystem}
%=============================================================================

\subsection{Central Timing Hub}

\begin{longtable}{p{5cm}p{8cm}}
\toprule
\textbf{Component} & \textbf{Specification} \\
\midrule
\endfirsthead
Master clock & Active hydrogen maser (MHM-2010 class) \\
Backup clock & Cesium beam standard (5071A class) \\
Optical reference & Optical frequency comb (Er-fiber, self-referenced) \\
Stability (1~s) & $\sigma_y(1\;\text{s}) < 2\times10^{-13}$ \\
Stability (1~day) & $\sigma_y(86{,}400\;\text{s}) < 5\times10^{-16}$ \\
UTC traceability & GNSS common-view + two-way satellite \\
\bottomrule
\end{longtable}

\subsection{Fiber Distribution Network}

Phase-coherent timing is distributed from the hub to each station via
stabilized optical fiber:

\begin{longtable}{p{5cm}p{8cm}}
\toprule
\textbf{Parameter} & \textbf{Specification} \\
\midrule
\endfirsthead
Fiber type & SMF-28e+ (ITU-T G.652.D) \\
Fiber length & 5.0~km per link (hub to each station) \\
Total fiber & 30~km (6 links, no ring) \\
Optical wavelength & 1542~nm (C-band, telecom compatible) \\
Stabilization method & Round-trip phase noise cancellation \\
Transfer stability & $\sigma_y(1\;\text{s}) < 10^{-15}$ \\
Transfer accuracy & $< 100$~ps (after calibration) \\
Phase noise floor & $< -120$~dBc/Hz at 1~Hz offset \\
Fiber route & Buried (frost-depth), armored conduit \\
Splice loss budget & $< 0.5$~dB total per link \\
\bottomrule
\end{longtable}

\subsubsection{Phase Stabilization}

Each fiber link uses an active noise cancellation system:
\begin{enumerate}
\item A narrow-linewidth CW laser at the hub is split: one arm sent
  through the fiber, the other retained as local reference.
\item At the far end (station), the signal is partially reflected back
  to the hub via a Faraday mirror.
\item The round-trip phase is compared to the local reference;
  the error signal drives an acousto-optic modulator (AOM) that
  pre-compensates the fiber phase noise.
\item The compensated optical signal at the station is converted to
  an RF reference (10~MHz, 100~MHz) via optical-to-electronic
  conversion.
\item The rotor phase-lock loop (PLL) locks to this reference.
\end{enumerate}

\subsection{Rotor Phase-Lock Protocol}

Each rotor is phase-locked to the distributed timing reference:
\begin{itemize}
\item An optical encoder on the rotor shaft provides $10^4$
  counts/revolution.
\item A digital PLL compares encoder phase to the station timing reference
  divided to the target rotation frequency.
\item Loop bandwidth: 10~Hz (sufficient to track rotor dynamics).
\item Phase-lock precision: $< 1~\mu$rad at the modulation frequency
  ($2\omega_{\rm rot}$).
\item Phase transients from disturbances (e.g., seismic) are flagged
  and the data segment is excluded in post-processing.
\end{itemize}

%=============================================================================
\section{Receiver Subsystem}
%=============================================================================

\subsection{Fixed Reference Receiver}

Located at the network center (hub), the fixed reference receiver
provides continuous monitoring of all six emitter signals and serves
as the primary performance benchmark.

\begin{longtable}{p{5cm}p{8cm}}
\toprule
\textbf{Component} & \textbf{Specification} \\
\midrule
\endfirsthead
$\psi$-sensor (primary) & Transportable optical lattice clock (Sr or Yb) \\
$\psi$-sensor (backup) & High-finesse optical cavity (ULE, 30~cm) \\
Gradient sensor & Superconducting gravity gradiometer (SGG) \\
$\psi$-stability & $\sigma_\psi < 10^{-17}$ at $\tau = 1$~s \\
Gradient sensitivity & $< 10^{-3}$~E ($10^{-12}$~s$^{-2}$) \\
Sampling rate & 100~Hz (all channels synchronous) \\
Data acquisition & 24-bit ADC, GPS-disciplined timing \\
Processing & Real-time spectral analysis + batch post-processing \\
\bottomrule
\end{longtable}

\subsection{Mobile Receiver Unit}

Transportable receiver for field positioning tests:

\begin{longtable}{p{5cm}p{8cm}}
\toprule
\textbf{Component} & \textbf{Specification} \\
\midrule
\endfirsthead
$\psi$-sensor & Compact optical clock (Ca+ or Sr+) or CSAC + cavity \\
Gradient sensor & MEMS accelerometer triad + tilt correction \\
Platform & Vehicle-mounted (van or truck) \\
Power & Vehicle alternator + battery bank, 3~kW \\
GNSS receiver & Multi-constellation RTK (for ground truth) \\
INS & Tactical-grade MEMS IMU \\
Data link & 4G/5G cellular + UHF radio backup \\
Processing & Embedded FPGA + ARM processor \\
\bottomrule
\end{longtable}

\subsection{Receiver Signal Processing}

\subsubsection{Spectral Separation}

The receiver separates contributions from individual emitters using:
\begin{enumerate}
\item \textbf{Bandpass filtering:} Digital IIR filters centered at each
  emitter modulation frequency ($f_k = 10.0, 12.5, \ldots, 22.5$~Hz)
  with bandwidth $\Delta f = 0.5$~Hz.
\item \textbf{Lock-in detection:} Phase-sensitive detection at each $f_k$,
  referenced to the known emitter modulation schedule.
\item \textbf{Coherent integration:} Accumulate in-phase and quadrature
  components over integration time $T_{\rm int}$. SNR improves as
  $\sqrt{T_{\rm int}/T_{\rm sample}}$.
\end{enumerate}

\subsubsection{Position Computation}

After source separation and coherent integration:
\begin{enumerate}
\item Extract $\psi_{E_k}^{\rm meas}$ (amplitude) and
  $\phi_{E_k}^{\rm meas}$ (phase) for each emitter $k$.
\item Compute range estimates $R_k = GM_{E_k}/(c^2\psi_{E_k}^{\rm meas})$.
\item Initialize position from prior (GNSS, INS, or previous $\psi$-fix).
\item Execute Gauss-Newton iteration (Sec.~5.5 of companion paper).
\item Apply curvature verification for integrity.
\item Output PVT solution with covariance and integrity flag.
\end{enumerate}

%=============================================================================
\section{Reference GNSS Subsystem}
%=============================================================================

The reference GNSS subsystem provides ground truth for validating
$\psi$-PNT performance.

\begin{longtable}{p{5cm}p{8cm}}
\toprule
\textbf{Component} & \textbf{Specification} \\
\midrule
\endfirsthead
Base station & Multi-constellation GNSS reference station at hub \\
Rover receiver & Matching GNSS receiver on mobile unit \\
RTK link & UHF radio or cellular NTRIP \\
RTK accuracy & $< 10$~mm horizontal, $< 20$~mm vertical \\
Post-processing & PPP-AR for $< 5$~mm accuracy \\
Timing & GNSS common-view for UTC comparison \\
Survey & All emitter positions surveyed to $< 1$~mm via static GNSS \\
\bottomrule
\end{longtable}

%=============================================================================
\section{Communications and Data Management}
%=============================================================================

\subsection{Station-to-Hub Communications}

\begin{longtable}{p{5cm}p{8cm}}
\toprule
\textbf{Parameter} & \textbf{Specification} \\
\midrule
\endfirsthead
Primary link & Dedicated fiber (wavelength-multiplexed with timing) \\
Data rate & 10~Mbps per station (continuous telemetry) \\
Backup link & 4G/5G cellular modem \\
Protocol & TCP/IP with custom telemetry protocol \\
Telemetry contents & Rotor state, environmental sensors, $\psi$-sensor data,
  health/status \\
Latency & $< 10$~ms (primary), $< 100$~ms (backup) \\
\bottomrule
\end{longtable}

\subsection{Data Storage and Processing}

\begin{longtable}{p{5cm}p{8cm}}
\toprule
\textbf{Parameter} & \textbf{Specification} \\
\midrule
\endfirsthead
Central server & Rack-mounted, dual Xeon, 512~GB RAM, 100~TB NAS \\
Data rate (total) & $\sim 60$~Mbps (6 stations + receivers) \\
Storage per day & $\sim 650$~GB raw data \\
Real-time processing & FPGA-accelerated spectral analysis \\
Batch processing & GPU-accelerated Bayesian parameter estimation \\
Data format & HDF5 with metadata and provenance tracking \\
Backup & Daily incremental to offsite NAS \\
\bottomrule
\end{longtable}

%=============================================================================
\section{Site Selection Criteria}
%=============================================================================

\subsection{Geophysical Requirements}

\begin{enumerate}
\item \textbf{Seismic quietude:} Site noise PSD below NLNM (New Low
  Noise Model) in the 10--25~Hz band. Avoid proximity to highways,
  railways, and industrial facilities.
\item \textbf{Geological stability:} Bedrock or consolidated sediment.
  Avoid active fault zones, karst terrain, and areas with significant
  subsidence.
\item \textbf{Topography:} Gently rolling or flat terrain preferred for
  uniform line-of-sight between stations. Maximum elevation variation
  $< 100$~m across network.
\item \textbf{Hydrology:} Low water table variability. Avoid floodplains
  and areas with large seasonal groundwater fluctuation.
\end{enumerate}

\subsection{Infrastructure Requirements}

\begin{enumerate}
\item \textbf{Power:} Grid electricity within 1~km of each station site.
\item \textbf{Fiber route:} Existing fiber duct or road shoulder
  available for burial.
\item \textbf{Road access:} All-weather road to each station for
  equipment delivery and maintenance.
\item \textbf{Land use:} Secure land tenure (lease or purchase) for
  $> 5$~year operational period.
\item \textbf{Cellular coverage:} 4G/5G coverage at all stations for
  backup communications.
\end{enumerate}

\subsection{Candidate Site Types}

\begin{itemize}
\item Rural agricultural land (low cultural noise, flat terrain, grid power);
\item University research farms (institutional support, existing infrastructure);
\item Military training areas (security, controlled access, low interference);
\item Decommissioned airfields (flat, surveyed, infrastructure available).
\end{itemize}

%=============================================================================
\section{Test Program}
%=============================================================================

\subsection{Phase 1: Single Emitter Characterization (Months 1--6)}

\paragraph{Objectives.}
\begin{enumerate}
\item Commission first emitter station (E1) and central hub.
\item Verify rotor operation: balance, speed stability, phase lock.
\item Deploy fixed reference receiver at hub.
\item Measure received $\psi$-signal (or upper limit) at multiple ranges.
\item Characterize sensor noise floor and systematic effects.
\end{enumerate}

\paragraph{Test matrix.}
\begin{center}
\begin{tabular}{cccc}
\toprule
\textbf{Range (m)} & \textbf{Integration (s)} & \textbf{Expected $\psi$} & \textbf{Measurement} \\
\midrule
100 & 3600 & Model-dependent & Spectral analysis \\
500 & 3600 & Model-dependent & Spectral analysis \\
1000 & 14400 & Model-dependent & Spectral analysis \\
2000 & 86400 & Model-dependent & Long integration \\
5000 & 604800 & Model-dependent & Week-long campaign \\
\bottomrule
\end{tabular}
\end{center}

\paragraph{Success criteria.}
Detection of coherent signal at the modulation frequency above the noise
floor at $\geq 3\sigma$ confidence at $\geq 1$ range, OR establishment
of upper bound on DFD coupling enhancement ratio $G_{\rm eff}/G$.

\subsection{Phase 2: Two-Emitter Baseline (Months 7--12)}

\paragraph{Objectives.}
\begin{enumerate}
\item Commission second emitter (E4, opposite side of network).
\item Demonstrate spectral separation of two emitter signals.
\item Measure differential $\psi$ along the baseline.
\item Characterize direction-of-arrival estimation from gradient.
\end{enumerate}

\paragraph{Success criteria.}
Successful spectral separation with $> 20$~dB isolation between
channels. Differential $\psi$ consistent with emitter geometry.

\subsection{Phase 3: Full Network Operation (Months 13--24)}

\paragraph{Objectives.}
\begin{enumerate}
\item Commission remaining four emitters (E2, E3, E5, E6).
\item Demonstrate simultaneous 6-emitter operation with spectral separation.
\item Compute first $\psi$-based position fix.
\item Characterize PDOP across the network.
\item Demonstrate curvature tensor verification.
\item Measure position accuracy vs.\ integration time.
\end{enumerate}

\paragraph{Success criteria.}
$\psi$-derived position fix consistent with GNSS/RTK ground truth within
error budget. Curvature tensor residual consistent with null spoofing
hypothesis.

\subsection{Phase 4: Hybrid GNSS+$\psi$ Fusion (Months 19--30)}

\paragraph{Objectives.}
\begin{enumerate}
\item Implement Kalman filter fusing GNSS and $\psi$ measurements.
\item Conduct simulated GPS spoofing attacks (GPS signal simulator).
\item Demonstrate spoofing detection via $\psi$-curvature residual.
\item Characterize detection probability vs.\ spoofing displacement.
\item Test holdover performance during simulated GNSS denial.
\end{enumerate}

\paragraph{Success criteria.}
Spoofing detection probability $> 99.9\%$ for displacements $> 10$~m.
Holdover position drift $< 10\times$ pure INS drift rate.

\subsection{Phase 5: Application Demonstrations (Months 25--36)}

\paragraph{Objectives.}
\begin{enumerate}
\item Autonomous vehicle demonstration: vehicle-mounted receiver
  navigating a predefined route within the network.
\item Timing demonstration: comparison of $\psi$-phase timing with
  GNSS and maser references.
\item Indoor demonstration: receiver in a building within the network,
  demonstrating positioning without GNSS.
\item Demonstration of continuous operation over $> 30$~days.
\end{enumerate}

%=============================================================================
\section{Risk Assessment and Mitigation}
%=============================================================================

\begin{longtable}{p{3.5cm}p{2cm}p{2cm}p{5.5cm}}
\toprule
\textbf{Risk} & \textbf{Likelihood} & \textbf{Impact} & \textbf{Mitigation} \\
\midrule
\endfirsthead
DFD coupling too weak for detection & Medium & Critical &
Design receiver for maximum sensitivity; near-field tests first; characterize upper bound scientifically \\
\addlinespace
Rotor vibration coupling to sensor & High & Major &
Vibration isolation; remote sensor placement; common-mode rejection \\
\addlinespace
Fiber timing instability & Low & Moderate &
Redundant stabilization; temperature-controlled duct; regular calibration \\
\addlinespace
Seismic noise contamination & Medium & Major &
Site selection; seismic veto; matched-filter processing \\
\addlinespace
Emitter mass uncertainty & Low & Minor &
Precision mass measurement pre-installation; in-situ calibration \\
\addlinespace
Environmental drift (thermal, pressure) & Medium & Moderate &
Active environmental control; correlation analysis; model subtraction \\
\addlinespace
Power interruption & Low & Minor &
UPS + generator; phase-lock re-acquisition protocol \\
\addlinespace
Fiber damage & Low & Moderate &
Buried armored cable; splice repair kit on-site; backup wireless link \\
\bottomrule
\end{longtable}

%=============================================================================
\section{Budget Estimate}
%=============================================================================

\begin{longtable}{p{6cm}p{3cm}p{3cm}}
\toprule
\textbf{Item} & \textbf{Unit cost} & \textbf{Total} \\
\midrule
\endfirsthead
\multicolumn{3}{l}{\textbf{Emitter stations ($\times 6$)}} \\
\quad Tungsten rotor assembly & \$200K & \$1,200K \\
\quad Magnetic bearings + vacuum & \$150K & \$900K \\
\quad Drive motor + PLL electronics & \$50K & \$300K \\
\quad Building + foundation & \$100K & \$600K \\
\quad Power system (UPS + generator) & \$40K & \$240K \\
\quad Environmental monitoring & \$30K & \$180K \\
\addlinespace
\multicolumn{3}{l}{\textbf{Timing distribution}} \\
\quad Hydrogen maser & \$300K & \$300K \\
\quad Optical frequency comb & \$200K & \$200K \\
\quad Stabilized fiber (30~km installed) & \$100/m & \$3,000K \\
\quad Phase stabilization electronics ($\times 6$) & \$50K & \$300K \\
\addlinespace
\multicolumn{3}{l}{\textbf{Receiver subsystem}} \\
\quad Transportable optical clock & \$1,000K & \$1,000K \\
\quad Superconducting gradiometer & \$500K & \$500K \\
\quad Mobile receiver unit ($\times 2$) & \$200K & \$400K \\
\quad Data acquisition system & \$100K & \$100K \\
\addlinespace
\multicolumn{3}{l}{\textbf{Reference GNSS}} \\
\quad RTK base station & \$20K & \$20K \\
\quad RTK rover ($\times 2$) & \$15K & \$30K \\
\addlinespace
\multicolumn{3}{l}{\textbf{Infrastructure}} \\
\quad Central hub building & --- & \$200K \\
\quad Communications (fiber multiplexing) & --- & \$150K \\
\quad Data server + storage & --- & \$100K \\
\addlinespace
\multicolumn{3}{l}{\textbf{Operations (3 years)}} \\
\quad Personnel (6 FTE $\times$ 3 yr) & \$120K/yr & \$2,160K \\
\quad Power + maintenance & \$50K/yr & \$150K \\
\quad Site lease & \$30K/yr & \$90K \\
\midrule
\textbf{Total} & & \textbf{\$12,120K} \\
\bottomrule
\end{longtable}

The estimated total project cost is approximately \textbf{\$12.1~million}
over a 3-year period, comparable to a medium-scale physics experiment
or a single GPS ground station modernization.

%=============================================================================
\section{Schedule}
%=============================================================================

\begin{longtable}{p{2cm}p{4cm}p{6cm}}
\toprule
\textbf{Month} & \textbf{Phase} & \textbf{Milestones} \\
\midrule
\endfirsthead
1--3 & Design \& procurement & Detailed design review; long-lead procurement \\
4--6 & Hub + E1 construction & Hub operational; E1 rotor installed \\
7--9 & Phase 1 testing & Single-emitter characterization complete \\
10--12 & E4 construction + Phase 2 & Two-emitter baseline demonstrated \\
13--18 & Full network build-out & All 6 stations operational \\
19--24 & Phase 3 testing & First $\psi$-position fix; PDOP mapping \\
25--30 & Phase 4 testing & GNSS+$\psi$ fusion; spoofing detection \\
31--36 & Phase 5 demonstrations & Application demos; final reporting \\
\bottomrule
\end{longtable}

%=============================================================================
\section{Intellectual Property and Regulatory}
%=============================================================================

\subsection{Patent Coverage}

This demonstration network falls under DFD Patent \#9: ``$\psi$-Field
Navigation, Positioning, and Timing Systems'' (Gary Thomas Alcock),
covering:
\begin{itemize}
\item Navigation via overlapping $\psi$ corridors from constellation of
  emitters;
\item Receiver computation of coordinates via delay surface intersection;
\item Clock-free timing from $\psi$ phase;
\item Anti-spoofing via curvature tensor verification;
\item Hybrid GNSS+$\psi$ fusion architecture.
\end{itemize}

\subsection{Regulatory Considerations}

\begin{itemize}
\item No RF emissions: the $\psi$-emitter produces no electromagnetic
  radiation, so no spectrum licensing is required.
\item Building permits for station structures (standard construction).
\item Environmental assessment for fiber burial (standard utility).
\item Occupational safety: rotating machinery (standard industrial).
\item The non-electromagnetic nature of $\psi$-PNT means it falls outside
  current radio navigation regulatory frameworks (ITU, FCC), potentially
  requiring novel regulatory engagement.
\end{itemize}

%=============================================================================
\section{Conclusions}
%=============================================================================

This engineering design provides a complete, buildable specification for
the first terrestrial demonstration of $\psi$-field positioning,
navigation, and timing. The design addresses all major engineering
subsystems: emitter mass assemblies with phase-locked rotation, precision
timing distribution via stabilized fiber, multi-sensor receivers, and a
reference GNSS system for ground-truth validation.

The phased test program is structured to provide scientific value at
every stage---from single-emitter signal characterization through full
network positioning demonstrations. Even a null result at Phase 1
(non-detection of DFD-enhanced $\psi$-coupling) would establish important
upper bounds on the DFD coupling enhancement, contributing to the
theoretical development of the framework.

The total project cost of $\sim$\$12M and 3-year timeline are modest
compared to typical navigation system development programs, making this
an accessible first step toward demonstrating the $\psi$-PNT concept.

\end{document}
